\chapter{Diseño del Proyecto}

\section{Stack Tecnológico}

\begin{longtable}{>{\bfseries}p{3cm} p{4cm} p{5.5cm}}
  \toprule
  \textbf{Capa} & \textbf{Tecnología} & \textbf{Justificación} \\
  \midrule
  \endhead
  Mobile y Web & React Native + Expo SDK 54 & Permite una misma base de código para Android y rutas web públicas. \\[3pt]
  Navegación & Expo Router 6 & Ruteo basado en archivos con grupos por rol y rutas públicas. \\[3pt]
  Lenguaje & TypeScript & Tipado estático para reducir errores durante el desarrollo. \\[3pt]
  Backend & Supabase & PostgreSQL, Auth, Storage, Realtime y funciones SQL administradas. \\[3pt]
  Base de datos & PostgreSQL & Modelo relacional, integridad referencial, constraints, triggers y RPC. \\[3pt]
  Autenticación & Supabase Auth + \texttt{profiles} & Sesiones seguras y perfiles propios con roles \texttt{admin}, \texttt{mesero} y \texttt{chef}. \\[3pt]
  Seguridad & Row Level Security & Restricción de acceso a datos según rol y usuario autenticado. \\[3pt]
  Almacenamiento & Supabase Storage & Imágenes del menú y logo del restaurante mediante URLs públicas. \\[3pt]
  Tiempo real & Supabase Realtime & Sincronización de órdenes, mesas, llamadas e inventario. \\[3pt]
  Reportes & Funciones PostgreSQL/RPC & Cálculos server-side para ventas, egresos, ranking y ventas por mesero. \\[3pt]
  UI & Componentes MD3 propios & Interfaz consistente para admin, mesero y chef. \\[3pt]
  Versiones & Git + GitHub & Historial de cambios, ramas auxiliares y respaldo del repositorio. \\
  \bottomrule
  \caption{Stack tecnológico real del sistema}
\end{longtable}

\section{Modelo de Datos}

\subsection{Entidades y Relaciones}

El sistema utiliza una base de datos PostgreSQL administrada por Supabase. El
modelo se organiza alrededor de restaurantes, usuarios/perfiles, mesas, menú,
órdenes, gastos, llamadas de mesa e inventario.

\begin{figure}[H]
    \centering
    \includegraphics[width=0.95\textwidth]{img/erd_database.png}
    \caption{Modelo entidad-relación de la base de datos del sistema iComanda}
    \label{fig:erd_database}
\end{figure}

\begin{notabox}
El archivo \texttt{report/erd\_mermaid.md} contiene el diagrama ER en formato
Mermaid para regenerar la imagen \texttt{img/erd\_database.png} cuando se actualice
el modelo.
\end{notabox}

\clearpage

\subsection{Esquema de Tablas Reales}

\noindent\textbf{\texttt{profiles}}

\begin{tabular}{lll}
  \toprule
  \textbf{Campo} & \textbf{Tipo} & \textbf{Descripción} \\
  \midrule
  \texttt{id} & uuid & PK, referencia a \texttt{auth.users} \\
  \texttt{nombre} & text & Nombre del usuario \\
  \texttt{email} & text & Correo electrónico \\
  \texttt{rol} & enum & \texttt{mesero | admin | chef} \\
  \texttt{activo} & boolean & Indica si el usuario puede ingresar \\
  \texttt{created\_at} & timestamptz & Fecha de creación \\
  \bottomrule
\end{tabular}

\vspace{8pt}
\noindent\textbf{\texttt{restaurants}}

\begin{tabular}{lll}
  \toprule
  \textbf{Campo} & \textbf{Tipo} & \textbf{Descripción} \\
  \midrule
  \texttt{id} & uuid & PK \\
  \texttt{nombre} & text & Nombre del restaurante \\
  \texttt{slug} & text UNIQUE & Identificador usado en URLs públicas \\
  \texttt{direccion} & text & Dirección del restaurante \\
  \texttt{telefono} & text & Teléfono de contacto \\
  \texttt{logo\_url} & text & Logo público del restaurante \\
  \texttt{owner\_id} & uuid FK & Administrador responsable \\
  \texttt{created\_at} & timestamptz & Fecha de creación \\
  \bottomrule
\end{tabular}

\vspace{8pt}
\noindent\textbf{\texttt{tables}}

\begin{tabular}{lll}
  \toprule
  \textbf{Campo} & \textbf{Tipo} & \textbf{Descripción} \\
  \midrule
  \texttt{id} & uuid & PK \\
  \texttt{restaurant\_id} & uuid FK & Restaurante al que pertenece \\
  \texttt{numero} & integer & Número de mesa \\
  \texttt{capacidad} & integer & Capacidad aproximada \\
  \texttt{activa} & boolean & Estado activo/inactivo \\
  \texttt{last\_cleared\_at} & timestamptz & Última limpieza/liberación de mesa \\
  \texttt{created\_at} & timestamptz & Fecha de creación \\
  \bottomrule
\end{tabular}

\vspace{8pt}
\noindent\textbf{\texttt{menu\_items}}

\begin{tabular}{lll}
  \toprule
  \textbf{Campo} & \textbf{Tipo} & \textbf{Descripción} \\
  \midrule
  \texttt{id} & uuid & PK \\
  \texttt{restaurant\_id} & uuid FK & Restaurante propietario \\
  \texttt{nombre} & text & Nombre del plato \\
  \texttt{descripcion} & text & Descripción del plato \\
  \texttt{precio} & numeric(10,2) & Precio actual \\
  \texttt{categoria} & text & Categoría del menú \\
  \texttt{imagen\_url} & text & Imagen del plato \\
  \texttt{disponible} & boolean & Visible/disponible para venta \\
  \texttt{agotado} & boolean & Marcado temporal de cocina \\
  \texttt{created\_at} & timestamptz & Fecha de creación \\
  \bottomrule
\end{tabular}

\vspace{8pt}
\noindent\textbf{\texttt{orders}}

\begin{tabular}{lll}
  \toprule
  \textbf{Campo} & \textbf{Tipo} & \textbf{Descripción} \\
  \midrule
  \texttt{id} & uuid & PK \\
  \texttt{restaurant\_id} & uuid FK & Restaurante asociado \\
  \texttt{table\_id} & uuid FK & Mesa asociada \\
  \texttt{mesero\_id} & uuid FK & Mesero que abrió la orden \\
  \texttt{estado} & enum & \texttt{activa | lista | entregada | cancelada} \\
  \texttt{metodo\_pago} & enum & \texttt{efectivo | qr | tarjeta} \\
  \texttt{pago\_confirmado} & boolean & Confirmación final del pago \\
  \texttt{paid\_at} & timestamptz & Fecha efectiva de pago \\
  \texttt{total} & numeric(10,2) & Total de la orden \\
  \texttt{created\_at} & timestamptz & Fecha de apertura \\
  \texttt{closed\_at} & timestamptz & Fecha de cierre \\
  \bottomrule
\end{tabular}

\vspace{8pt}
\noindent\textbf{\texttt{order\_items}}

\begin{tabular}{lll}
  \toprule
  \textbf{Campo} & \textbf{Tipo} & \textbf{Descripción} \\
  \midrule
  \texttt{id} & uuid & PK \\
  \texttt{order\_id} & uuid FK & Orden padre \\
  \texttt{menu\_item\_id} & uuid FK & Plato original, puede quedar nulo \\
  \texttt{nombre} & text & Snapshot del nombre del plato \\
  \texttt{precio\_unitario} & numeric(10,2) & Snapshot del precio \\
  \texttt{cantidad} & integer & Cantidad solicitada \\
  \texttt{created\_at} & timestamptz & Fecha de registro \\
  \bottomrule
\end{tabular}

\vspace{8pt}
\noindent\textbf{\texttt{expenses}}

\begin{tabular}{lll}
  \toprule
  \textbf{Campo} & \textbf{Tipo} & \textbf{Descripción} \\
  \midrule
  \texttt{id} & uuid & PK \\
  \texttt{restaurant\_id} & uuid FK & Restaurante asociado \\
  \texttt{descripcion} & text & Descripción del gasto \\
  \texttt{monto} & numeric(10,2) & Monto del gasto \\
  \texttt{fecha} & date & Fecha contable \\
  \texttt{created\_by} & uuid FK & Usuario que lo registró \\
  \texttt{created\_at} & timestamptz & Fecha de creación \\
  \bottomrule
\end{tabular}

\vspace{8pt}
\noindent\textbf{\texttt{table\_calls}}

\begin{tabular}{lll}
  \toprule
  \textbf{Campo} & \textbf{Tipo} & \textbf{Descripción} \\
  \midrule
  \texttt{id} & uuid & PK \\
  \texttt{table\_id} & uuid FK & Mesa que solicita atención \\
  \texttt{restaurant\_id} & uuid FK & Restaurante asociado \\
  \texttt{atendida} & boolean & Indica si fue atendida \\
  \texttt{atendida\_at} & timestamptz & Momento de atención \\
  \texttt{created\_at} & timestamptz & Momento de llamada \\
  \bottomrule
\end{tabular}

\vspace{8pt}
\noindent\textbf{\texttt{materials}}

\begin{tabular}{lll}
  \toprule
  \textbf{Campo} & \textbf{Tipo} & \textbf{Descripción} \\
  \midrule
  \texttt{id} & uuid & PK \\
  \texttt{restaurant\_id} & uuid FK & Restaurante asociado \\
  \texttt{nombre} & text & Nombre del material o insumo \\
  \texttt{unidad} & text & Unidad de medida \\
  \texttt{stock\_minimo} & numeric & Umbral mínimo \\
  \texttt{activo} & boolean & Estado activo/inactivo \\
  \texttt{created\_at} & timestamptz & Fecha de creación \\
  \bottomrule
\end{tabular}

\vspace{8pt}
\noindent\textbf{\texttt{material\_stock\_movements}}

\begin{tabular}{lll}
  \toprule
  \textbf{Campo} & \textbf{Tipo} & \textbf{Descripción} \\
  \midrule
  \texttt{id} & uuid & PK \\
  \texttt{material\_id} & uuid FK & Material afectado \\
  \texttt{restaurant\_id} & uuid FK & Restaurante asociado \\
  \texttt{movement\_type} & enum & \texttt{consumo | reposicion} \\
  \texttt{quantity} & numeric & Cantidad movida \\
  \texttt{reason} & text & Motivo del movimiento \\
  \texttt{notes} & text & Notas adicionales \\
  \texttt{created\_by} & uuid FK & Usuario que registró el movimiento \\
  \texttt{created\_at} & timestamptz & Fecha del movimiento \\
  \bottomrule
\end{tabular}

\section{Estructura del Proyecto}

La estructura real del repositorio sigue las convenciones de Expo Router:

\begin{Verbatim}[fontsize=\small, frame=single, framesep=3mm]
iComanda/
|-- app/
|   |-- login.tsx
|   |-- _layout.tsx
|   |-- (admin)/       index, usuarios, mesas, restaurante, menu,
|   |                  pagos-qr, gastos, inventario, reports
|   |-- (mesero)/      index, menu, pedidos, pedido-nuevo, pedido/[id]
|   |-- (chef)/        index
|   |-- menu/[slug].tsx
|   \-- table/[slug]/[tableId].tsx
|-- components/
|   |-- admin/
|   |-- customer/
|   |-- md3/
|   \-- tables/
|-- context/auth.tsx
|-- lib/
|   |-- supabase.ts
|   |-- supabase-public.ts
|   |-- database.types.ts
|   |-- admin.ts
|   |-- order-receipt.ts
|   \-- hooks/useTables.ts
|-- supabase/
|   |-- migrations/
|   \-- functions/crear-usuario/
|-- report/
|-- app.json
|-- package.json
\-- tsconfig.json
\end{Verbatim}

\section{Roles y Accesos}

\begin{longtable}{p{6cm} ccc}
  \toprule
  \textbf{Funcionalidad} & \textbf{Mesero} & \textbf{Admin} & \textbf{Chef} \\
  \midrule
  \endhead
  Ver menú disponible & \checkmark & \checkmark & \checkmark \\
  Crear pedidos & \checkmark & --- & --- \\
  Ver pedidos activos & \checkmark & \checkmark & \checkmark \\
  Marcar pedido como listo & --- & --- & \checkmark \\
  Marcar pedido entregado/cancelado & \checkmark & \checkmark & --- \\
  Vista de mesas en tiempo real & \checkmark & --- & --- \\
  Atender llamada de cliente & \checkmark & --- & --- \\
  Gestionar menú & --- & \checkmark & --- \\
  Marcar plato agotado & --- & --- & \checkmark \\
  Gestionar mesas y restaurante & --- & \checkmark & --- \\
  Gestionar usuarios & --- & \checkmark & --- \\
  Confirmar pagos QR & --- & \checkmark & --- \\
  Registrar gastos & --- & \checkmark & --- \\
  Gestionar inventario & --- & \checkmark & --- \\
  Ver reportes & --- & \checkmark & --- \\
  \bottomrule
  \caption{Matriz real de roles y accesos del sistema}
\end{longtable}

\section{Rutas Públicas del Cliente}

Además de los roles autenticados, el sistema expone rutas públicas para clientes:

\begin{itemize}[leftmargin=1.5em]
  \item \texttt{/menu/[slug]}: carta digital pública del restaurante. Ejemplo local validado: \texttt{/menu/jimmy}.
  \item \texttt{/table/[slug]/[tableId]}: vista pública de mesa para ver menú y llamar al mesero. Ejemplo local validado: \texttt{/table/jimmy/b3383250-5517-4772-8bda-549bf27d44e1}.
\end{itemize}

Estas rutas no requieren una cuenta del cliente y están protegidas por políticas
RLS y consultas limitadas al restaurante/mesa correspondiente.
